\chapter{Graph Theory — Karger Min Cut Trace}
%%% generated by the blueprint pipeline from TCSlib.GraphTheory.KargerMinCutTrace
%%%%%%%%%%%%%%%%%%%%%%%%%%%%%%%%%%%%%%%%%%%%%%%%%%%%%%%%%%%%%%%%%%%%%%%%%%%%%%%
\section{Overview}

This chapter documents the algorithmic trace layer for Karger's contraction algorithm:
adaptive contraction runs of a multigraph together with a distinguished cut carried along
the contractions, a finite edge-occurrence sampling model, and the resulting survival
probability lower bound $2/(n(n-1))$ for a fixed min-cut.  It builds on the local
contraction facts (cut size preservation, min-cut preservation, degree and handshaking
bounds) proved for multigraphs elsewhere, and packages them both recursively and as a
telescoping product along an indexed trace.

%%%%%%%%%%%%%%%%%%%%%%%%%%%%%%%%%%%%%%%%%%%%%%%%%%%%%%%%%%%%%%%%%%%%%%%%%%%%%%%
\section{Declarations}

\begin{definition}[State of the contraction process]
\lean{KargerTrace.State}
\label{KargerTrace.State}
\leanok
\uses{MulCut, Multigraph}
A state consists of a multigraph $G$ over $\alpha$ together with a distinguished cut $C$
of $G$.  It records the current graph of the algorithm and the cut whose survival is
being tracked.
\end{definition}

\begin{definition}[Contraction of a state along a non-crossing edge]
\lean{KargerTrace.contractState}
\label{KargerTrace.contractState}
\leanok
\uses{KargerTrace.State, MulCut.contractedCut, MulCut.crossingEdges, Multigraph.contract}
Given a state $s$, an edge $e$ of $s.G$ that does not cross the distinguished cut $s.C$,
the contracted state has graph $s.G$ contracted along $e$ and cut the corresponding
contracted cut.
\end{definition}

\begin{definition}[Surviving step]
\lean{KargerTrace.SurvivingStep}
\label{KargerTrace.SurvivingStep}
\leanok
\uses{KargerTrace.State, KargerTrace.contractState, MulCut.crossingEdges}
The relation asserting that state $t$ arises from state $s$ by contracting some edge $e$
of $s.G$ that does not cross the distinguished cut of $s$.
\end{definition}

\begin{definition}[Surviving run]
\lean{KargerTrace.SurvivingRun}
\label{KargerTrace.SurvivingRun}
\leanok
\uses{KargerTrace.State, KargerTrace.SurvivingStep}
The inductively defined relation of $n$ consecutive surviving steps leading from a state
$s$ to a state $t$: the empty run has $n = 0$ and $t = s$, and a run of length $n+1$ is a
run of length $n$ followed by one surviving step.
\end{definition}

\begin{lemma}[Positivity of the edge count when an edge exists]
\lean{KargerTrace.edgeCount_pos_of_mem}
\label{KargerTrace.edgeCount_pos_of_mem}
\leanok
\uses{Multigraph, Multigraph.edgeCount}
\difficulty{1}
Let $G$ be a multigraph on a vertex type $\alpha$, and suppose $e$ is an edge of $G$.
Then the edge count of $G$, the number of edges counted with multiplicity, is strictly
positive.
\end{lemma}

\begin{lemma}[Contraction along a non-crossing edge preserves cut size]
\lean{KargerTrace.survivingStep_cut_size_eq}
\label{KargerTrace.survivingStep_cut_size_eq}
\leanok
\uses{KargerTrace.State, KargerTrace.SurvivingStep, MulCut.size, cut_size_preserved}
\difficulty{2}
Let $s$ and $t$ be states of the contraction process, each consisting of a multigraph
together with a distinguished cut of that multigraph. If $t$ arises from $s$ by a
surviving step — that is, $t$ is obtained by contracting some edge $e$ of the multigraph
of $s$ that does not cross the distinguished cut of $s$ — then the two distinguished
cuts have the same size; equivalently, the number of edges crossing the distinguished
cut of $t$, counted with multiplicity, equals the number of edges crossing the
distinguished cut of $s$.
\end{lemma}

\begin{lemma}[A surviving contraction preserves minimum-cut status]
\lean{KargerTrace.survivingStep_mincut}
\label{KargerTrace.survivingStep_mincut}
\leanok
\uses{IsMulMinCut, KargerTrace.State, KargerTrace.SurvivingStep, mincut_preserved_of_non_crossing}
\difficulty{2}
Let $s$ be a state of the contraction process, consisting of a multigraph $G$ together
with a distinguished cut $C$ of $G$, and let $t$ be a state obtained from $s$ by a
surviving step: that is, $t$ arises by contracting some edge $e$ of $G$ that does not
cross $C$, so that $t$ has graph $G/e$ and its distinguished cut is the cut induced by
$C$ on $G/e$. If $C$ is a minimum cut of $G$, then the distinguished cut of $t$ is a
minimum cut of $G/e$.
\end{lemma}

\begin{lemma}[A surviving contraction step removes one vertex]
\lean{KargerTrace.survivingStep_vertexCount_eq}
\label{KargerTrace.survivingStep_vertexCount_eq}
\leanok
\uses{KargerTrace.State, KargerTrace.SurvivingStep, KargerTrace.contractState, Multigraph.vertexCount, contract_vertex_count}
\difficulty{3}
Let $s$ and $t$ be states of the contraction process, each a multigraph together with a
distinguished cut, and suppose that $t$ arises from $s$ by a surviving step, that is,
$t$ is obtained by contracting some edge $e$ of the multigraph of $s$ that does not
cross the distinguished cut of $s$. Then the number of vertices of the multigraph of $t$
equals the number of vertices of the multigraph of $s$ minus $1$ (with the difference
taken in $\bbn$).
\end{lemma}

\begin{lemma}[Positive edge count under a surviving step]
\lean{KargerTrace.survivingStep_edgeCount_pos}
\label{KargerTrace.survivingStep_edgeCount_pos}
\leanok
\uses{KargerTrace.State, KargerTrace.SurvivingStep, KargerTrace.edgeCount_pos_of_mem, Multigraph.edgeCount}
\difficulty{2}
Let $s$ and $t$ be states of the contraction process, each a multigraph equipped with a
distinguished cut. If $t$ arises from $s$ by a surviving step — that is, by contracting
some edge of the multigraph of $s$ that does not cross the distinguished cut of $s$ —
then the multigraph of $s$ has positive edge count.
\end{lemma}

\begin{lemma}[Invariance of cut size along a surviving run]
\lean{KargerTrace.survivingRun_cut_size_eq}
\label{KargerTrace.survivingRun_cut_size_eq}
\leanok
\uses{KargerTrace.State, KargerTrace.SurvivingRun, KargerTrace.survivingStep_cut_size_eq, MulCut.size}
\difficulty{2}
Let $s$ and $t$ be states of the contraction process, each consisting of a multigraph
together with a distinguished cut, and suppose that $t$ is obtained from $s$ by a
surviving run of length $n$, that is, by $n$ consecutive contractions of edges that do
not cross the distinguished cut of the current state. Then the distinguished cut of $t$
has the same size as the distinguished cut of $s$, where the size of a cut is the number
of its crossing edges counted with multiplicity.
\end{lemma}

\begin{lemma}[Surviving runs preserve minimality of the tracked cut]
\lean{KargerTrace.survivingRun_mincut}
\label{KargerTrace.survivingRun_mincut}
\leanok
\uses{IsMulMinCut, KargerTrace.State, KargerTrace.SurvivingRun, KargerTrace.survivingStep_mincut}
\difficulty{2}
Let $s$ and $t$ be states of the contraction process, each consisting of a multigraph
together with a distinguished cut of that multigraph, and suppose that $t$ is obtained
from $s$ by a surviving run of some length $n$, that is, by $n$ consecutive contractions
of edges that never cross the distinguished cut. If the distinguished cut of $s$ is a
minimum cut of its multigraph, then the distinguished cut of $t$ is a minimum cut of its
multigraph.
\end{lemma}

\begin{lemma}[Vertex count after an $n$-step surviving run]
\lean{KargerTrace.survivingRun_vertexCount_eq}
\label{KargerTrace.survivingRun_vertexCount_eq}
\leanok
\uses{KargerTrace.State, KargerTrace.SurvivingRun, KargerTrace.survivingStep_vertexCount_eq, Multigraph.vertexCount}
\difficulty{3}
Let $s$ and $t$ be states of the contraction process, each consisting of a multigraph
together with a distinguished cut, and suppose that $t$ is reached from $s$ by a
surviving run of length $n$, that is, by $n$ consecutive contractions of edges that do
not cross the tracked cut. Writing $G_s$ and $G_t$ for the multigraphs of $s$ and $t$,
and $\lvert V(G)\rvert$ for the number of vertices of a multigraph $G$, we have \[
\lvert V(G_t)\rvert = \lvert V(G_s)\rvert - n, \] where the subtraction is truncated
subtraction on $\bbn$ (equal to $0$ once $n$ exceeds $\lvert V(G_s)\rvert$).
\end{lemma}

\begin{lemma}[Karger min-cut edge-count invariant along a surviving run]
\lean{KargerTrace.survivingRun_edge_count_invariant}
\label{KargerTrace.survivingRun_edge_count_invariant}
\leanok
\uses{IsMulMinCut, KargerTrace.State, KargerTrace.SurvivingRun, KargerTrace.survivingRun_cut_size_eq, KargerTrace.survivingRun_mincut, MulCut.size, Multigraph.edgeCount, Multigraph.vertexCount, edge_count_lower_bound, mincut_le_degree}
\difficulty{4}
Let $s$ and $t$ be states of the contraction process, each consisting of a multigraph
together with a distinguished cut, and suppose that $t$ is obtained from $s$ by a
surviving run of some length $n$, that is, by $n$ successive contractions of edges that
do not cross the distinguished cut. If the distinguished cut $C$ of $s$ is a minimum cut
of the multigraph $s.G$, with size $k$ (its number of crossing edges), then the number
of vertices $|V|$ and the number of edges $|E|$ of the final multigraph $t.G$ satisfy
\[
  k \cdot |V| \;\le\; 2\,|E|.
\]
\end{lemma}

\begin{lemma}[Indexed edge-count invariant for surviving runs]
\lean{KargerTrace.survivingRun_edge_count_invariant_indexed}
\label{KargerTrace.survivingRun_edge_count_invariant_indexed}
\leanok
\uses{IsMulMinCut, KargerTrace.State, KargerTrace.SurvivingRun, KargerTrace.survivingRun_edge_count_invariant, KargerTrace.survivingRun_vertexCount_eq, MulCut.size, Multigraph.edgeCount, Multigraph.vertexCount}
\difficulty{2}
Let $s$ and $t$ be states of the contraction process, each consisting of a multigraph
together with a distinguished cut, and suppose $t$ is reached from $s$ by a surviving
run of $n$ steps. Write $C$ for the distinguished cut of $s$ and $G$ for its multigraph,
and let $G'$ denote the multigraph of $t$. If $C$ is a minimum cut of $G$, then
\[
  \abs{C}\,\bigl(\abs{V(G)} - n\bigr) \;\le\; 2\,\abs{E(G')},
\]
where $\abs{C}$ is the size of the cut (its number of crossing edges), $\abs{V(G)}$ is
the number of vertices of $G$, $\abs{E(G')}$ is the number of edges of $G'$ counted with
multiplicity, and the difference $\abs{V(G)} - n$ is truncated at $0$.
\end{lemma}

\begin{lemma}[Algebraic one-step survival bound]
\lean{KargerTrace.survivalProb_ge_factor}
\label{KargerTrace.survivalProb_ge_factor}
\leanok
\uses{survivalProb}
\difficulty{4}
Let $c$, $n$, and $m$ be natural numbers with $n \ge 2$, $m > 0$, and $cn \le 2m$. Then
the one-step survival probability $1 - \frac{c}{m}$ satisfies
\[
  \frac{n-2}{n} \;\le\; 1 - \frac{c}{m}.
\]
\end{lemma}

\begin{lemma}[One-step survival bound for a minimum cut]
\lean{KargerTrace.survivingStep_survivalProb_ge_factor}
\label{KargerTrace.survivingStep_survivalProb_ge_factor}
\leanok
\uses{IsMulMinCut, KargerTrace.State, KargerTrace.SurvivingStep, KargerTrace.survivalProb_ge_factor, KargerTrace.survivingStep_edgeCount_pos, MulCut.size, Multigraph.edgeCount, Multigraph.vertexCount, edge_count_lower_bound, mincut_le_degree, survivalProb}
\difficulty{3}
Let $s$ be a state of the contraction process, consisting of a multigraph $G$ together
with a distinguished cut $C$ of $G$. Suppose that a surviving step is possible from $s$,
meaning that $G$ has an edge that does not cross $C$, whose contraction — merging the
endpoints of that edge, deleting the resulting self-loops, and passing to the cut that
$C$ induces on the contracted graph — yields a further state of the process. Assume
moreover that $C$ is a minimum cut of $G$ and that $G$ has at least two vertices.
Writing $n$ for the number of vertices of $G$, $k$ for the size of $C$ (its number of
crossing edges, counted with multiplicity), and $m$ for the number of edges of $G$
(counted with multiplicity), one has
\[
  \frac{n-2}{n} \;\le\; 1 - \frac{k}{m},
\]
so that $\frac{n-2}{n}$ is at most the one-step survival probability of $C$.
\end{lemma}

\begin{lemma}[One-step survival lower bound in edge contraction]
\lean{KargerTrace.survivingRun_survivalProb_ge_factor}
\label{KargerTrace.survivingRun_survivalProb_ge_factor}
\leanok
\uses{IsMulMinCut, KargerTrace.State, KargerTrace.SurvivingRun, KargerTrace.SurvivingStep, KargerTrace.survivalProb_ge_factor, KargerTrace.survivingRun_cut_size_eq, KargerTrace.survivingRun_mincut, KargerTrace.survivingStep_edgeCount_pos, MulCut.size, Multigraph.edgeCount, Multigraph.vertexCount, edge_count_lower_bound, mincut_le_degree, survivalProb}
\difficulty{4}
Let $s$, $t$, and $u$ be states of the contraction process, each a multigraph paired
with a distinguished cut. Suppose $t$ is reached from $s$ by a surviving run of $i$
steps, that a further surviving step leads from $t$ to $u$, that the distinguished cut
of $s$ is a minimum cut of the multigraph of $s$, and that the multigraph of $t$ has at
least two vertices. Writing $n_t$ for the number of vertices of the multigraph of $t$,
$c$ for the size of the distinguished cut of $s$, and $m$ for the number of edges of the
multigraph of $t$, then
\[
  \frac{n_t - 2}{n_t} \;\le\; 1 - \frac{c}{m}.
\]
\end{lemma}

\begin{definition}[Non-crossing edges of a state]
\lean{KargerTrace.nonCrossingEdges}
\label{KargerTrace.nonCrossingEdges}
\leanok
\uses{Crosses, KargerTrace.State}
The sub-multiset of edges of $s.G$ that do not cross the distinguished cut $s.C$; these
are exactly the edge choices that keep the cut alive.
\end{definition}

\begin{definition}[Number of non-crossing choices]
\lean{KargerTrace.nonCrossingChoiceCount}
\label{KargerTrace.nonCrossingChoiceCount}
\leanok
\uses{KargerTrace.State, KargerTrace.nonCrossingEdges}
The cardinality (with multiplicity) of the multiset of non-crossing edges of a state.
\end{definition}

\begin{definition}[Uniform one-step survival probability]
\lean{KargerTrace.uniformSurvivalProb}
\label{KargerTrace.uniformSurvivalProb}
\leanok
\uses{KargerTrace.State, KargerTrace.nonCrossingChoiceCount, Multigraph.edgeCount}
The real number obtained as the number of non-crossing edge choices divided by the total
edge count of $s.G$, i.e.\ the probability that a uniformly chosen edge occurrence does not
cross the cut.
\end{definition}

\begin{lemma}[Non-crossing edge count equals edges minus cut size]
\lean{KargerTrace.nonCrossingChoiceCount_eq_edgeCount_sub_cutSize}
\label{KargerTrace.nonCrossingChoiceCount_eq_edgeCount_sub_cutSize}
\leanok
\uses{Crosses, KargerTrace.State, KargerTrace.nonCrossingChoiceCount, MulCut.size, Multigraph.edgeCount, crossingEdges_eq_filter_crosses}
\difficulty{4}
Let $s$ be a state of the contraction process, consisting of a multigraph $G$ together
with a distinguished cut $C$, whose two sides are given by a vertex subset $S$ and its
complement. Then the number of edges of $G$ that do not cross $C$ — those having no
endpoint in $S$ or no endpoint outside $S$, counted with multiplicity — equals the total
edge count of $G$ minus the size of $C$, the difference being taken as truncated
subtraction in $\bbn$:
\[
\#\{\text{non-crossing edges of }s\} \;=\; \abs{E(G)} - \abs{C}.
\]
\end{lemma}

\begin{lemma}[Cut size bounded by edge count]
\lean{KargerTrace.cut_size_le_edgeCount}
\label{KargerTrace.cut_size_le_edgeCount}
\leanok
\uses{KargerTrace.State, MulCut.crossingEdges, MulCut.size, Multigraph.edgeCount}
\difficulty{2}
Let $G$ be a multigraph and let $C$ be a cut of $G$, given by a vertex subset $S$ that
is nonempty and has nonempty complement. Then the size of $C$ — the number of edges of
$G$ having exactly one endpoint in $S$, counted with multiplicity — is at most the total
number of edges of $G$, again counted with multiplicity.
\end{lemma}

\begin{lemma}[Uniform edge choice realizes the survival probability]
\lean{KargerTrace.uniformSurvivalProb_eq_survivalProb}
\label{KargerTrace.uniformSurvivalProb_eq_survivalProb}
\leanok
\uses{KargerTrace.State, KargerTrace.cut_size_le_edgeCount, KargerTrace.nonCrossingChoiceCount_eq_edgeCount_sub_cutSize, KargerTrace.uniformSurvivalProb, MulCut.size, Multigraph.edgeCount, survivalProb}
\difficulty{3}
Let $s$ be a state of the contraction process, consisting of a multigraph $G$ together
with a distinguished cut $C$, and suppose $G$ has at least one edge, so that its edge
count $m$ is positive. Then the uniform one-step survival probability of $s$ — the
number of non-crossing edge choices divided by $m$ — equals
\[
  1 - \frac{c}{m},
\]
where $c$ is the size of the cut $C$, that is, its number of crossing edges (each
counted with multiplicity).
\end{lemma}

\begin{definition}[Sample space of edge occurrences]
\lean{KargerTrace.edgeChoiceSpace}
\label{KargerTrace.edgeChoiceSpace}
\leanok
\uses{KargerTrace.State}
The finite set of pairs (edge, occurrence index) obtained from the edge multiset of $s.G$;
the second coordinate distinguishes parallel copies of the same edge, so that a uniform
choice on this set is a uniform choice of an edge occurrence.
\end{definition}

\begin{definition}[Surviving edge occurrence]
\lean{KargerTrace.edgeChoiceSurvives}
\label{KargerTrace.edgeChoiceSurvives}
\leanok
\uses{Crosses, KargerTrace.State}
The predicate on an edge occurrence saying that its underlying edge does not cross the
distinguished cut of $s$.
\end{definition}

\begin{lemma}[Cardinality of the edge-occurrence space]
\lean{KargerTrace.edgeChoiceSpace_card}
\label{KargerTrace.edgeChoiceSpace_card}
\leanok
\uses{KargerTrace.State, KargerTrace.edgeChoiceSpace, Multigraph.edgeCount}
\difficulty{1}
Let $s$ be a state of the contraction process, consisting of a multigraph $G$ together
with a cut of $G$, and let its sample space of edge occurrences be the finite set of
pairs (edge, occurrence index) drawn from the edge multiset of $G$. Then this sample
space has cardinality equal to the edge count of $G$, that is, the number of edges of
$G$ counted with multiplicity.
\end{lemma}

\begin{lemma}[Filtering enumerated multiset occurrences by element]
\lean{KargerTrace.toEnumFinset_filter_fst_card_eq_filter_card}
\label{KargerTrace.toEnumFinset_filter_fst_card_eq_filter_card}
\leanok
\difficulty{4}
Let $\alpha$ be a type with decidable equality, let $m$ be a finite multiset over
$\alpha$, and let $p$ be a decidable predicate on $\alpha$. Form the enumeration set of
$m$, consisting of all pairs $(a, i)$ with $a$ an element of $m$ and $0 \le i <
\text{count}_m(a)$, so that each occurrence of each element of $m$ is represented
exactly once. Then the number of such pairs $(a,i)$ whose first coordinate $a$ satisfies
$p$ equals the size (counted with multiplicity) of the sub-multiset of $m$ consisting of
the elements satisfying $p$.
\end{lemma}

\begin{lemma}[Surviving edge occurrences equal non-crossing edge count]
\lean{KargerTrace.edgeChoiceSurvivalCount_eq_nonCrossingChoiceCount}
\label{KargerTrace.edgeChoiceSurvivalCount_eq_nonCrossingChoiceCount}
\leanok
\uses{Crosses, KargerTrace.State, KargerTrace.edgeChoiceSpace, KargerTrace.edgeChoiceSurvives, KargerTrace.nonCrossingChoiceCount, KargerTrace.nonCrossingEdges, KargerTrace.toEnumFinset_filter_fst_card_eq_filter_card}
\difficulty{1}
Let $s$ be a state of the contraction process, consisting of a loopless multigraph $G$
together with a distinguished cut $S \subseteq V(G)$, and enumerate the edge multiset of
$G$ as a finite set of edge occurrences $(e, i)$, where the index $i$ distinguishes
parallel copies of the same edge. Say that an occurrence $(e, i)$ survives when its
underlying edge $e$ does not cross $S$, that is, when it is not the case that one
endpoint of $e$ lies in $S$ and the other lies outside $S$. Then the number of surviving
edge occurrences equals the number of edges of $G$ that do not cross $S$, counted with
multiplicity.
\end{lemma}

\begin{definition}[Survival ratio of the occurrence space]
\lean{KargerTrace.edgeChoiceSurvivalRatio}
\label{KargerTrace.edgeChoiceSurvivalRatio}
\leanok
\uses{KargerTrace.State, KargerTrace.edgeChoiceSpace, KargerTrace.edgeChoiceSurvives}
The ratio of the number of surviving edge occurrences to the total number of edge
occurrences, i.e.\ the concrete one-step survival probability of the finite sampling model.
\end{definition}

\begin{lemma}[Survival ratio equals the uniform survival probability]
\lean{KargerTrace.edgeChoiceSurvivalRatio_eq_uniformSurvivalProb}
\label{KargerTrace.edgeChoiceSurvivalRatio_eq_uniformSurvivalProb}
\leanok
\uses{KargerTrace.State, KargerTrace.edgeChoiceSpace_card, KargerTrace.edgeChoiceSurvivalCount_eq_nonCrossingChoiceCount, KargerTrace.edgeChoiceSurvivalRatio, KargerTrace.uniformSurvivalProb}
\difficulty{1}
Let $s$ be a state of the contraction process, consisting of a multigraph $G$ together
with a distinguished cut $C$ of $G$. Form the survival ratio by enumerating the edge
occurrences of $G$ — the pairs of an edge with an index distinguishing its parallel
copies — and dividing the number of occurrences whose underlying edge does not cross $C$
by the total number of occurrences. Then this ratio equals the uniform survival
probability of $s$, namely the number of edges of $G$ that do not cross $C$, counted
with multiplicity, divided by the total edge count of $G$.
\end{lemma}

\begin{lemma}[Edge component of a sampled occurrence lies in the graph]
\lean{KargerTrace.edgeChoice_mem_edges}
\label{KargerTrace.edgeChoice_mem_edges}
\leanok
\uses{KargerTrace.State, KargerTrace.edgeChoiceSpace}
\difficulty{2}
Let $s$ be a state of the contraction process, with underlying multigraph $G$, and let
$(e, k)$ be any element of the sample space of edge occurrences of $s$ — a pair
consisting of an edge together with an index distinguishing parallel copies. Then the
edge $e$ is a member of the edge multiset of $G$.
\end{lemma}

\begin{lemma}[Surviving edge occurrences do not cross the cut]
\lean{KargerTrace.not_mem_crossingEdges_of_edgeChoiceSurvives}
\label{KargerTrace.not_mem_crossingEdges_of_edgeChoiceSurvives}
\leanok
\uses{Crosses, KargerTrace.State, KargerTrace.edgeChoiceSpace, KargerTrace.edgeChoiceSurvives, MulCut.crossingEdges, crossingEdges_eq_filter_crosses}
\difficulty{3}
Let $s = (G, C)$ be a state of the contraction process, consisting of a multigraph $G$
and a distinguished cut $C$ with side $S$, and let $(e, i)$ be an edge occurrence in the
sample space of $s$, so that $e$ is an edge of $G$ and $i$ indexes its parallel copies.
If this occurrence survives — that is, if $e$ does not cross $S$, meaning it is not the
case that $e$ has one endpoint in $S$ and another endpoint outside $S$ — then $e$ does
not belong to the multiset of crossing edges of $C$.
\end{lemma}

\begin{definition}[Successor state of a surviving occurrence]
\lean{KargerTrace.stateAfterSurvivingChoice}
\label{KargerTrace.stateAfterSurvivingChoice}
\leanok
\uses{KargerTrace.State, KargerTrace.contractState, KargerTrace.edgeChoiceSpace, KargerTrace.edgeChoiceSurvives, KargerTrace.edgeChoice_mem_edges, KargerTrace.not_mem_crossingEdges_of_edgeChoiceSurvives}
Given a state $s$ and a sampled edge occurrence that does not cross the distinguished cut,
this is the state obtained by contracting $s$ along the underlying edge.
\end{definition}

\begin{lemma}[Sampled surviving occurrence yields a surviving step]
\lean{KargerTrace.stateAfterSurvivingChoice_step}
\label{KargerTrace.stateAfterSurvivingChoice_step}
\leanok
\uses{KargerTrace.State, KargerTrace.SurvivingStep, KargerTrace.edgeChoiceSpace, KargerTrace.edgeChoiceSurvives, KargerTrace.edgeChoice_mem_edges, KargerTrace.not_mem_crossingEdges_of_edgeChoiceSurvives, KargerTrace.stateAfterSurvivingChoice}
\difficulty{1}
Let $s$ be a state of the contraction process, consisting of a multigraph $G$ together
with a distinguished cut $C$ of $G$, and let a sampled edge occurrence be given in the
sample space of $s$ — a pair (edge, occurrence index) drawn from the edge multiset of
$G$ — whose underlying edge does not cross $C$. Then the pair formed by $s$ and the
successor state obtained by contracting $s$ along that edge occurrence stands in the
surviving-step relation; that is, the successor state arises from $s$ by contracting
some edge of $G$ that does not cross $C$.
\end{lemma}

\begin{definition}[Recursive survival probability]
\lean{KargerTrace.survivalEventProb}
\label{KargerTrace.survivalEventProb}
\leanok
\uses{KargerTrace.State, KargerTrace.edgeChoiceSpace, KargerTrace.edgeChoiceSurvives, KargerTrace.stateAfterSurvivingChoice}
The finite probability that the distinguished cut survives the next given number of random
edge-occurrence choices, defined by recursion on the number of steps: zero steps have
probability $1$, and one further step averages over all edge occurrences, contributing $0$
for occurrences that cross the cut and the recursive value at the contracted state
otherwise.
\end{definition}

\begin{lemma}[Zero-step survival probability equals one]
\lean{KargerTrace.survivalEventProb_zero}
\label{KargerTrace.survivalEventProb_zero}
\leanok
\uses{KargerTrace.State, KargerTrace.survivalEventProb}
\difficulty{1}
For every state $s$ — a multigraph $G$ together with a distinguished cut $C$ of $G$ —
the recursive survival probability of the distinguished cut over zero random
edge-occurrence choices equals $1$.
\end{lemma}

\begin{lemma}[One-step recursion for the survival probability]
\lean{KargerTrace.survivalEventProb_succ}
\label{KargerTrace.survivalEventProb_succ}
\leanok
\uses{KargerTrace.State, KargerTrace.edgeChoiceSpace, KargerTrace.edgeChoiceSurvives, KargerTrace.stateAfterSurvivingChoice, KargerTrace.survivalEventProb}
\difficulty{1}
Let $s$ be a state, consisting of a multigraph $G$ together with a distinguished cut $C$
of $G$, and let $\mathrm{steps}$ be a natural number. Write $P_k(s)$ for the survival
probability of the cut over $k$ random edge-occurrence choices starting from $s$. Then
\[
P_{\mathrm{steps}+1}(s)
  \;=\;
  \frac{1}{N}\sum_{\text{occurrences } c}
  \begin{cases}
P_{\mathrm{steps}}\!\bigl(s_c\bigr) & \text{if the edge of } c \text{ does not cross }
C,\\[4pt]
    0 & \text{otherwise,}
  \end{cases}
\]
where the sum ranges over the finite set of edge occurrences of $G$ (each parallel copy
of an edge counted separately), $N$ is the number of such occurrences, and $s_c$ denotes
the state obtained by contracting $s$ along the underlying edge of a surviving
occurrence $c$.
\end{lemma}

\begin{lemma}[Survivor count over the attached sample space]
\lean{KargerTrace.edgeChoice_attach_filter_card}
\label{KargerTrace.edgeChoice_attach_filter_card}
\leanok
\uses{KargerTrace.State, KargerTrace.edgeChoiceSpace, KargerTrace.edgeChoiceSurvives}
\difficulty{2}
Let $s$ be a state of the contraction process, consisting of a multigraph $G$ together
with a distinguished cut $S$ of $G$, and let $\Omega$ be its sample space of edge
occurrences — the finite set of pairs $(e, i)$ where $e$ is an edge of $G$ and $i$ is an
occurrence index distinguishing parallel copies of that edge. Call an occurrence
*surviving* when its underlying edge $e$ does not cross $S$, that is, when it is not the
case that one endpoint of $e$ lies in $S$ and the other lies outside $S$. Then the
number of surviving occurrences is the same whether they are counted directly among the
elements of $\Omega$ or among the elements of the equinumerous copy of $\Omega$ obtained
by regarding each member together with the fact that it belongs to $\Omega$; in both
counts the survival test is applied to the underlying occurrence, and the two
cardinalities agree.
\end{lemma}

\begin{lemma}[Counting lower bound for a conditional sum]
\lean{KargerTrace.card_mul_le_sum_ite_of_lower_bound}
\label{KargerTrace.card_mul_le_sum_ite_of_lower_bound}
\leanok
\difficulty{4}
Let $S$ be a finite set, let $p$ be a predicate on $S$, and let $f\colon S \to \bbr$.
Suppose $B \in \bbr$ is a lower bound for $f$ on the elements satisfying $p$, that is,
$B \le f(x)$ for every $x \in S$ with $p(x)$. Then
\[
\abs{\{x \in S : p(x)\}} \cdot B \;\le\; \sum_{x \in S} \begin{cases} f(x) & \text{if }
p(x),\\ 0 & \text{otherwise.}\end{cases}
\]
\end{lemma}

\begin{lemma}[One-step lower bound for the survival probability]
\lean{KargerTrace.survivalEventProb_step_lower_bound}
\label{KargerTrace.survivalEventProb_step_lower_bound}
\leanok
\uses{KargerTrace.State, KargerTrace.card_mul_le_sum_ite_of_lower_bound, KargerTrace.edgeChoiceSpace, KargerTrace.edgeChoiceSurvivalRatio, KargerTrace.edgeChoiceSurvives, KargerTrace.edgeChoice_attach_filter_card, KargerTrace.stateAfterSurvivingChoice, KargerTrace.survivalEventProb, KargerTrace.survivalEventProb_succ}
\difficulty{5}
Let $s$ be a state of the contraction process — a multigraph together with a
distinguished cut — and assume that its sample space of edge occurrences is nonempty.
Fix a natural number $k$ and a real number $B$, and suppose that for every edge
occurrence that survives, meaning its underlying edge does not cross the distinguished
cut, the recursive survival probability over $k$ further choices, evaluated at the
successor state obtained by contracting $s$ along that occurrence, is at least $B$. Then
the recursive survival probability over $k+1$ choices at $s$ is at least the survival
ratio of $s$ multiplied by $B$.
\end{lemma}

\begin{lemma}[Survival ratio bounded below by the Karger factor]
\lean{KargerTrace.edgeChoiceSurvivalRatio_ge_factor}
\label{KargerTrace.edgeChoiceSurvivalRatio_ge_factor}
\leanok
\uses{IsMulMinCut, KargerTrace.State, KargerTrace.cut_size_le_edgeCount, KargerTrace.edgeChoiceSurvivalRatio, KargerTrace.edgeChoiceSurvivalRatio_eq_uniformSurvivalProb, KargerTrace.survivalProb_ge_factor, KargerTrace.uniformSurvivalProb_eq_survivalProb, MulCut.size, Multigraph.edgeCount, Multigraph.vertexCount, edge_count_lower_bound, mincut_le_degree}
\difficulty{3}
Let $s$ be a state of the contraction process — a multigraph $G$ together with a
distinguished cut $C$ of $G$ — and write $n$ for the number of vertices of $G$. Suppose
that $C$ is a minimum cut of $G$, that the size of $C$ is strictly positive, and that $n
\ge 2$. Then the survival ratio of the edge-occurrence sample space of $s$ is at least
$\tfrac{n-2}{n}$; that is,
\[
  \frac{n-2}{n} \;\le\; \rho(s),
\]
where $\rho(s)$ is the fraction of edge occurrences whose underlying edge does not cross
$C$.
\end{lemma}

\begin{theorem}[Karger contraction survival lower bound]
\lean{KargerTrace.survivalEventProb_lower_bound}
\label{KargerTrace.survivalEventProb_lower_bound}
\leanok
\uses{IsMulMinCut, KargerTrace.State, KargerTrace.SurvivingStep, KargerTrace.cut_size_le_edgeCount, KargerTrace.edgeChoiceSpace, KargerTrace.edgeChoiceSpace_card, KargerTrace.edgeChoiceSurvivalRatio, KargerTrace.edgeChoiceSurvivalRatio_ge_factor, KargerTrace.stateAfterSurvivingChoice, KargerTrace.stateAfterSurvivingChoice_step, KargerTrace.survivalEventProb, KargerTrace.survivalEventProb_step_lower_bound, KargerTrace.survivingStep_cut_size_eq, KargerTrace.survivingStep_mincut, KargerTrace.survivingStep_vertexCount_eq, MulCut.size, Multigraph.edgeCount, Multigraph.vertexCount}
\difficulty{6}
Let $s$ be a state of the contraction process, consisting of a multigraph $G$ together
with a distinguished cut $C$, and let $n$ be a natural number with $2 \le n$. Suppose
that $C$ is a minimum cut of $G$, that $C$ has positive size, and that $G$ has exactly
$n$ vertices. Then the recursive survival probability that $C$ survives $n-2$ successive
uniformly random edge-occurrence contractions is at least
\[
  \frac{2}{n(n-1)}.
\]
\end{theorem}

\begin{theorem}[Karger survival lower bound at the initial state]
\lean{KargerTrace.survivalEventProb_initial_lower_bound}
\label{KargerTrace.survivalEventProb_initial_lower_bound}
\leanok
\uses{IsMulMinCut, KargerTrace.State, KargerTrace.survivalEventProb, KargerTrace.survivalEventProb_lower_bound, MulCut.size, Multigraph.vertexCount}
\difficulty{1}
Let $s$ be a state of the contraction process, consisting of a multigraph $G$ together
with a distinguished cut $C$ of $G$, and write $n$ for the number of vertices of $G$.
Suppose that $C$ is a minimum cut of $G$, that its size is positive, and that $n \ge 2$.
Then the probability that $C$ survives $n-2$ successive uniformly random edge-occurrence
contractions is at least
\[
  \frac{2}{n(n-1)}.
\]
\end{theorem}

\begin{theorem}[Karger min-cut survival probability lower bound]
\lean{KargerTrace.fixed_mincut_survival_probability_lower_bound}
\label{KargerTrace.fixed_mincut_survival_probability_lower_bound}
\leanok
\uses{IsMulMinCut, KargerTrace.State, KargerTrace.survivalEventProb, KargerTrace.survivalEventProb_initial_lower_bound, MulCut.size, Multigraph.vertexCount}
\difficulty{1}
Let $s$ be a state consisting of a multigraph $G$ on $n = \abs{V(G)}$ vertices together
with a distinguished cut $C$ of $G$, and suppose that $C$ is a minimum cut of $G$, that
its size is positive, and that $n \ge 2$. Then the probability that $C$ survives $n-2$
successive uniformly random contractions of edge occurrences is at least
\[
\frac{2}{n(n-1)}.
\]
\end{theorem}

\begin{definition}[Graph-only occurrence space]
\lean{KargerTrace.graphChoiceSpace}
\label{KargerTrace.graphChoiceSpace}
\leanok
\uses{Multigraph}
The finite set of edge occurrences of a multigraph $G$, the same construction as the
state-level sample space but without reference to a distinguished cut.
\end{definition}

\begin{lemma}[Agreement of the graph-only and state occurrence spaces]
\lean{KargerTrace.graphChoiceSpace_eq_edgeChoiceSpace}
\label{KargerTrace.graphChoiceSpace_eq_edgeChoiceSpace}
\leanok
\uses{KargerTrace.State, KargerTrace.edgeChoiceSpace, KargerTrace.graphChoiceSpace}
\difficulty{1}
Let $s$ be a state of the contraction process, consisting of a multigraph $G$ together
with a distinguished cut of $G$. Then the graph-only occurrence space built from $G$
coincides, as a finite set of edge occurrences $(e, k) \in \mathrm{Sym}_2\,\alpha \times
\bbn$, with the state's sample space of edge occurrences; that is, both are the set of
edge occurrences of the state's underlying multigraph.
\end{lemma}

\begin{lemma}[Edge component of an occurrence lies in the edge multiset]
\lean{KargerTrace.graphChoice_mem_edges}
\label{KargerTrace.graphChoice_mem_edges}
\leanok
\uses{KargerTrace.graphChoiceSpace, Multigraph}
\difficulty{2}
Let $G$ be a multigraph, and let its graph-only occurrence space be the finite set of
edge occurrences of $G$, whose elements are pairs $(e, k)$ consisting of an edge $e$
together with an index $k$ distinguishing the copies of $e$ among the parallel edges of
$G$. Then for every occurrence $(e, k)$ in this space, the edge $e$ belongs to the edge
multiset of $G$.
\end{lemma}

\begin{definition}[Graph after a sampled contraction]
\lean{KargerTrace.graphAfterChoice}
\label{KargerTrace.graphAfterChoice}
\leanok
\uses{KargerTrace.graphChoiceSpace, KargerTrace.graphChoice_mem_edges, Multigraph, Multigraph.contract}
The multigraph obtained by contracting $G$ along the edge underlying a sampled edge
occurrence.
\end{definition}

\begin{lemma}[Sampled contraction removes exactly one vertex]
\lean{KargerTrace.graphAfterChoice_vertexCount_eq}
\label{KargerTrace.graphAfterChoice_vertexCount_eq}
\leanok
\uses{KargerTrace.graphAfterChoice, KargerTrace.graphChoiceSpace, KargerTrace.graphChoice_mem_edges, Multigraph, Multigraph.vertexCount, contract_vertex_count}
\difficulty{2}
Let $G$ be a multigraph on a vertex type $\alpha$ with decidable equality, and let $c$
be a sampled edge occurrence of $G$, i.e.\ an element of its finite set of edge
occurrences. Writing $G_c$ for the multigraph obtained by contracting $G$ along the edge
underlying $c$, the number of vertices of $G_c$ equals the number of vertices of $G$
minus one.
\end{lemma}

\begin{lemma}[State contraction agrees with graph contraction]
\lean{KargerTrace.stateAfterSurvivingChoice_graph_eq}
\label{KargerTrace.stateAfterSurvivingChoice_graph_eq}
\leanok
\uses{KargerTrace.State, KargerTrace.edgeChoiceSpace, KargerTrace.edgeChoiceSurvives, KargerTrace.graphAfterChoice, KargerTrace.graphChoiceSpace, KargerTrace.graphChoiceSpace_eq_edgeChoiceSpace, KargerTrace.stateAfterSurvivingChoice}
\difficulty{1}
Let $s$ be a state of the contraction process, comprising a multigraph $G$ together with
a distinguished cut $C$, and let $(e, i)$ be an edge occurrence of $G$ drawn from its
sample space of edge occurrences, whose underlying edge $e$ does not cross $C$. Then the
multigraph carried by the successor state obtained by contracting $s$ along this
occurrence is equal to the multigraph obtained by contracting $G$ itself along the same
occurrence, namely the contraction $G/e$.
\end{lemma}

\begin{definition}[Graph-only Karger run]
\lean{KargerTrace.KargerGraphRun}
\label{KargerTrace.KargerGraphRun}
\leanok
\uses{KargerTrace.graphAfterChoice, KargerTrace.graphChoiceSpace, Multigraph}
The inductive relation asserting that $H$ is reachable from $G$ by a given number of
successive arbitrary edge-occurrence contractions.  It models the actual graph evolution
of the algorithm without tracking any distinguished cut.
\end{definition}

\begin{lemma}[Vertex count after a graph-only Karger run]
\lean{KargerTrace.kargerGraphRun_vertexCount_eq}
\label{KargerTrace.kargerGraphRun_vertexCount_eq}
\leanok
\uses{KargerTrace.KargerGraphRun, KargerTrace.graphAfterChoice_vertexCount_eq, Multigraph, Multigraph.vertexCount}
\difficulty{3}
Let $G$ and $H$ be multigraphs on a vertex type $\alpha$, and let $\mathrm{steps} \in
\bbn$. If $H$ arises from $G$ by a graph-only Karger run consisting of $\mathrm{steps}$
successive edge-occurrence contractions, then the number of vertices of $H$ equals the
number of vertices of $G$ minus $\mathrm{steps}$, using truncated subtraction of natural
numbers. Thus each contraction removes exactly one vertex, and after $\mathrm{steps}$ of
them the vertex count has dropped by $\mathrm{steps}$.
\end{lemma}

\begin{lemma}[A full Karger run terminates with two vertices]
\lean{KargerTrace.kargerGraphRun_full_vertexCount_eq_two}
\label{KargerTrace.kargerGraphRun_full_vertexCount_eq_two}
\leanok
\uses{KargerTrace.KargerGraphRun, KargerTrace.kargerGraphRun_vertexCount_eq, Multigraph, Multigraph.vertexCount}
\difficulty{2}
Let $G$ and $H$ be multigraphs on a common vertex type, and suppose $G$ has at least two
vertices. If $H$ is obtained from $G$ by a graph-only Karger run consisting of exactly
$n - 2$ successive edge-occurrence contractions, where $n$ is the number of vertices of
$G$, then $H$ has exactly two vertices.
\end{lemma}

\begin{lemma}[Termination of a full Karger run]
\lean{KargerTrace.kargerGraphRun_full_output}
\label{KargerTrace.kargerGraphRun_full_output}
\leanok
\uses{KargerOutput, KargerTrace.KargerGraphRun, KargerTrace.kargerGraphRun_full_vertexCount_eq_two, Multigraph, Multigraph.vertexCount}
\difficulty{1}
Let $G$ be a multigraph with at least two vertices, write $n$ for its vertex count, and
suppose $H$ is obtained from $G$ by a graph-only Karger run consisting of exactly $n -
2$ successive edge-occurrence contractions. Then $H$ satisfies the termination condition
of the algorithm; that is, $H$ has exactly two vertices.
\end{lemma}

\begin{lemma}[The zeroth occurrence of an edge lies in the occurrence space]
\lean{KargerTrace.choice_mem_graphChoiceSpace_of_edge_mem}
\label{KargerTrace.choice_mem_graphChoiceSpace_of_edge_mem}
\leanok
\uses{KargerTrace.graphChoiceSpace, Multigraph}
\difficulty{2}
Let $G$ be a multigraph on a vertex type $\alpha$, and let $e$ be an edge of $G$, that
is, a member of the edge multiset of $G$. Then the occurrence $(e, 0)$ belongs to the
graph-only occurrence space of $G$; in other words, the pair consisting of $e$ together
with the index $0$ is one of the enumerated edge occurrences of $G$.
\end{lemma}

\begin{lemma}[Surviving step realized by a sampled edge occurrence]
\lean{KargerTrace.survivingStep_graphChoice}
\label{KargerTrace.survivingStep_graphChoice}
\leanok
\uses{KargerTrace.State, KargerTrace.SurvivingStep, KargerTrace.choice_mem_graphChoiceSpace_of_edge_mem, KargerTrace.graphAfterChoice, KargerTrace.graphChoiceSpace}
\difficulty{3}
Let $s$ and $t$ be states of the contraction process, each consisting of a multigraph
together with a distinguished cut of it. If $t$ arises from $s$ by a surviving step—that
is, by contracting some edge of $s$'s multigraph that does not cross its distinguished
cut—then there is an edge occurrence in the graph-only occurrence space of $s$'s
multigraph whose sampled contraction produces exactly the multigraph of $t$.
\end{lemma}

\begin{lemma}[Surviving runs project to graph runs]
\lean{KargerTrace.survivingRun_graphRun}
\label{KargerTrace.survivingRun_graphRun}
\leanok
\uses{KargerTrace.KargerGraphRun, KargerTrace.State, KargerTrace.SurvivingRun, KargerTrace.survivingStep_graphChoice}
\difficulty{3}
Let $s$ and $t$ be states of the contraction process over a vertex type $\alpha$, each
consisting of a multigraph together with a distinguished cut of it. If $t$ is reached
from $s$ by a surviving run of some number $\mathrm{steps}$ of consecutive surviving
steps — each contracting an edge that does not cross the current distinguished cut —
then the underlying multigraph of $t$ is reached from the underlying multigraph of $s$
by a graph-only Karger run of the same length $\mathrm{steps}$, that is, by
$\mathrm{steps}$ successive edge-occurrence contractions.
\end{lemma}

\begin{lemma}[Termination of a full surviving run at two vertices]
\lean{KargerTrace.survivingRun_full_output}
\label{KargerTrace.survivingRun_full_output}
\leanok
\uses{KargerOutput, KargerTrace.State, KargerTrace.SurvivingRun, KargerTrace.kargerGraphRun_full_output, KargerTrace.survivingRun_graphRun, Multigraph.vertexCount}
\difficulty{1}
Let $s$ be a state of the contraction process, consisting of a multigraph $G$ together
with a distinguished cut of $G$, and write $n$ for the number of vertices of $G$;
suppose $n \ge 2$. If a state $t$ is reached from $s$ by a surviving run of exactly $n -
2$ consecutive surviving steps — each step contracting an edge of the current graph that
does not cross the distinguished cut — then the multigraph of $t$ satisfies the
termination condition of the algorithm, that is, it has exactly two vertices.
\end{lemma}

\begin{definition}[Adaptive graph trace]
\lean{KargerTrace.KargerGraphTrace}
\label{KargerTrace.KargerGraphTrace}
\leanok
\uses{KargerTrace.graphAfterChoice, KargerTrace.graphChoiceSpace, Multigraph}
An explicit indexed execution of the contraction loop: a family of graphs indexed by
$\mathbb{N}$, an edge occurrence of the graph at each index below $\text{steps}$, and the
requirement that the next graph is the contraction of the current one along that
occurrence.  The choice type depends on the current graph, matching the adaptive nature of
the process.
\end{definition}

\begin{lemma}[Prefixes of a graph trace are graph runs]
\lean{KargerTrace.kargerGraphTrace_prefix_run}
\label{KargerTrace.kargerGraphTrace_prefix_run}
\leanok
\uses{KargerTrace.KargerGraphRun, KargerTrace.KargerGraphTrace, KargerTrace.graphAfterChoice}
\difficulty{3}
Let $T$ be an adaptive graph trace of length $\text{steps}$, so that $T$ records a graph
at each index, a sampled edge occurrence at each index below $\text{steps}$, and
produces each successive graph as the contraction of its predecessor along that
occurrence. Then for every index $i \le \text{steps}$, the graph-only Karger run
relation holds for $i$ steps between the initial graph $T_0$ and the graph $T_i$; that
is, $T_i$ is reachable from $T_0$ by exactly $i$ successive edge-occurrence
contractions.
\end{lemma}

\begin{lemma}[Vertex count decreases by one per contraction step]
\lean{KargerTrace.kargerGraphTrace_vertexCount_eq}
\label{KargerTrace.kargerGraphTrace_vertexCount_eq}
\leanok
\uses{KargerTrace.KargerGraphTrace, KargerTrace.kargerGraphRun_vertexCount_eq, KargerTrace.kargerGraphTrace_prefix_run, Multigraph.vertexCount}
\difficulty{1}
Let $(G_i)_{i \in \bbn}$ be an adaptive contraction trace over $\text{steps}$ rounds on
a vertex type with decidable equality: a family of multigraphs in which, for each index
$i < \text{steps}$, the graph $G_{i+1}$ is obtained from $G_i$ by contracting the edge
underlying a sampled edge occurrence of $G_i$. Then for every $i \le \text{steps}$, the
number of vertices of $G_i$ equals the initial vertex count minus $i$; that is, $\lvert
V(G_i)\rvert = \lvert V(G_0)\rvert - i$.
\end{lemma}

\begin{lemma}[Karger's algorithm terminates with two vertices]
\lean{KargerTrace.kargerGraphTrace_full_output}
\label{KargerTrace.kargerGraphTrace_full_output}
\leanok
\uses{KargerOutput, KargerTrace.KargerGraphTrace, KargerTrace.kargerGraphTrace_vertexCount_eq, Multigraph.vertexCount}
\difficulty{2}
Let $n \ge 2$, and let $T$ be an adaptive graph trace of $n-2$ steps whose initial
multigraph $T_0$ has exactly $n$ vertices. Then the multigraph $T_{n-2}$ reached after
all $n-2$ contraction steps is a valid output of the algorithm; that is, it has exactly
two vertices.
\end{lemma}

\begin{theorem}[Karger contraction terminates at two vertices]
\lean{KargerTrace.adaptive_karger_trace_outputs_two_vertices}
\label{KargerTrace.adaptive_karger_trace_outputs_two_vertices}
\leanok
\uses{KargerOutput, KargerTrace.KargerGraphTrace, KargerTrace.kargerGraphTrace_full_output, Multigraph.vertexCount}
\difficulty{1}
Let $n \ge 2$, and let $T$ be an adaptive graph trace of $n-2$ contraction steps whose
initial multigraph $T_0$ has exactly $n$ vertices. Then the final multigraph $T_{n-2}$,
obtained after the $n-2$ successive edge contractions prescribed by the trace, satisfies
the termination condition: it has exactly two vertices.
\end{theorem}

\begin{theorem}[Karger contraction: min-cut survival bound and termination]
\lean{KargerTrace.karger_algorithmic_trace_summary}
\label{KargerTrace.karger_algorithmic_trace_summary}
\leanok
\uses{IsMulMinCut, KargerOutput, KargerTrace.KargerGraphTrace, KargerTrace.State, KargerTrace.kargerGraphTrace_full_output, KargerTrace.survivalEventProb, KargerTrace.survivalEventProb_initial_lower_bound, MulCut.size, Multigraph.vertexCount}
\difficulty{2}
Let $s$ be a state consisting of a multigraph $G$ with $n$ vertices together with a
distinguished cut $C$ of $G$, and suppose that $C$ is a minimum cut, that its size is
strictly positive, and that $n \ge 2$. Then two conclusions hold simultaneously. First,
the recursive survival probability that $C$ is preserved through the next $n-2$ random
edge-occurrence contractions of $s$ is at least
\[
\frac{2}{n(n-1)}.
\]
Second, every adaptive graph trace of $n-2$ contraction steps whose initial graph is $G$
terminates in the output condition of the algorithm, that is, its graph after $n-2$
steps has exactly two vertices.
\end{theorem}

\begin{theorem}[Karger contraction: minimum-cut survival bound and termination]
\lean{KargerTrace.adaptive_karger_trace_main}
\label{KargerTrace.adaptive_karger_trace_main}
\leanok
\uses{IsMulMinCut, KargerOutput, KargerTrace.KargerGraphTrace, KargerTrace.State, KargerTrace.karger_algorithmic_trace_summary, KargerTrace.survivalEventProb, MulCut.size, Multigraph.vertexCount}
\difficulty{1}
Let $s$ be a state of the contraction process, consisting of a multigraph $G$ together
with a distinguished cut $C$ of $G$, and write $n$ for the number of vertices of $G$.
Suppose that $C$ is a minimum cut of $G$, that $C$ has positive size, and that $n \ge
2$. Let $p$ denote the probability that $C$ survives $n-2$ successive uniformly random
edge-occurrence contractions. Then
\[
\frac{2}{n(n-1)} \;\le\; p,
\]
and, moreover, every adaptive graph trace of $n-2$ contraction steps whose initial graph
is $G$ ends in a multigraph with exactly two vertices.
\end{theorem}

\begin{definition}[Indexed surviving trace]
\lean{KargerTrace.SurvivingTrace}
\label{KargerTrace.SurvivingTrace}
\leanok
\uses{KargerTrace.State, KargerTrace.SurvivingStep}
A family of states indexed by $\mathbb{N}$ such that consecutive states at indices below
$\text{steps}$ are related by a surviving step.  This is the product-friendly counterpart
of a surviving run: it exposes the state present at each index.
\end{definition}

\begin{definition}[Graph trace of a surviving trace]
\lean{KargerTrace.survivingTrace_toGraphTrace}
\label{KargerTrace.survivingTrace_toGraphTrace}
\leanok
\uses{KargerTrace.KargerGraphTrace, KargerTrace.SurvivingTrace, KargerTrace.survivingStep_graphChoice}
The adaptive graph trace obtained from an indexed surviving trace by forgetting the
distinguished cuts and choosing, at each step, an occurrence realizing the surviving step.
\end{definition}

\begin{lemma}[Prefixes of a surviving trace are surviving runs]
\lean{KargerTrace.survivingTrace_prefix_run}
\label{KargerTrace.survivingTrace_prefix_run}
\leanok
\uses{KargerTrace.SurvivingRun, KargerTrace.SurvivingTrace}
\difficulty{3}
Let $T$ be a surviving trace of length $\text{steps}$ over a vertex type $\alpha$, that
is, a family of states $(T_j)_{j \in \bbn}$ in which each consecutive pair $T_j,
T_{j+1}$ with $j < \text{steps}$ is related by a surviving step. Then for every index
$i$ with $i \le \text{steps}$, the states $T_0, T_1, \dots, T_i$ constitute a surviving
run of length $i$ from the initial state $T_0$ to the state $T_i$.
\end{lemma}

\begin{lemma}[Vertex count along a surviving trace]
\lean{KargerTrace.survivingTrace_vertexCount_eq}
\label{KargerTrace.survivingTrace_vertexCount_eq}
\leanok
\uses{KargerTrace.SurvivingTrace, KargerTrace.survivingRun_vertexCount_eq, KargerTrace.survivingTrace_prefix_run, Multigraph.vertexCount}
\difficulty{1}
Let $T$ be a surviving trace of length $\mathrm{steps}$ over a vertex type $\alpha$: a
family of states $(G_j, C_j)$ indexed by $j \in \bbn$, in which each consecutive pair
below $\mathrm{steps}$ is joined by a surviving step (the graph $G_{j+1}$ arises from
$G_j$ by contracting an edge that does not cross the distinguished cut). Then for every
index $i \le \mathrm{steps}$, the number of vertices of the graph at index $i$ equals
the number of vertices of the initial graph at index $0$ minus $i$; that is, $\lvert
V(G_i)\rvert = \lvert V(G_0)\rvert - i$.
\end{lemma}

\begin{lemma}[Per-step survival bound along a surviving trace]
\lean{KargerTrace.survivingTrace_survivalProb_ge_factor}
\label{KargerTrace.survivingTrace_survivalProb_ge_factor}
\leanok
\uses{IsMulMinCut, KargerTrace.SurvivingTrace, KargerTrace.survivingRun_survivalProb_ge_factor, KargerTrace.survivingTrace_prefix_run, MulCut.size, Multigraph.edgeCount, Multigraph.vertexCount, survivalProb}
\difficulty{1}
Let $T = (G_0, C_0), (G_1, C_1), \dots$ be a surviving trace: a family of states indexed
by $\bbn$ in which each state at an index below $\text{steps}$ passes to its successor
by contracting some edge of its current multigraph that does not cross its distinguished
cut. Suppose the initial cut $C_0$ is a minimum cut of $G_0$, and fix an index $i <
\text{steps}$ at which the multigraph $G_i$ has at least two vertices. Writing $n_i$ for
the number of vertices of $G_i$, $m_i$ for the number of edges of $G_i$ (counted with
multiplicity), and $c$ for the size of the initial cut $C_0$ (its number of crossing
edges), one has
\[
  \frac{n_i - 2}{n_i} \;\le\; 1 - \frac{c}{m_i}.
\]
\end{lemma}

\begin{lemma}[Karger step factors bound the survival-probability product]
\lean{KargerTrace.survivingTrace_factor_product_le_survival_product}
\label{KargerTrace.survivingTrace_factor_product_le_survival_product}
\leanok
\uses{IsMulMinCut, KargerTrace.SurvivingTrace, KargerTrace.survivingTrace_survivalProb_ge_factor, MulCut.size, Multigraph.edgeCount, Multigraph.vertexCount, survivalProb}
\difficulty{3}
Let $T$ be a surviving trace of $\mathrm{steps}$ steps, with states $s_0, s_1, \dots$;
write $G_i$ for the multigraph and $C_i$ for the distinguished cut in state $s_i$, and
let $n_i$ and $m_i$ denote the vertex count and edge count of $G_i$. Suppose the initial
cut $C_0$ is a minimum cut of $G_0$, and that $n_i \ge 2$ for every index $i <
\mathrm{steps}$. Then, writing $c$ for the size of $C_0$,
\[
  \prod_{i < \mathrm{steps}} \frac{n_i - 2}{n_i}
  \;\le\;
  \prod_{i < \mathrm{steps}} \left(1 - \frac{c}{m_i}\right),
\]
the right-hand factors being the one-step survival probabilities for cut size $c$ and
edge count $m_i$.
\end{lemma}

\begin{lemma}[Telescoping form of the surviving-trace factor product]
\lean{KargerTrace.survivingTrace_factor_product_eq_telescope}
\label{KargerTrace.survivingTrace_factor_product_eq_telescope}
\leanok
\uses{KargerTrace.SurvivingTrace, KargerTrace.survivingTrace_vertexCount_eq, Multigraph.vertexCount}
\difficulty{4}
Let $T$ be a surviving trace of length $\text{steps}$, and for each index $i$ write
$v_i$ for the number of vertices of the multigraph in its state at index $i$. Suppose
that $\text{steps} = n - 2$ and that the initial graph has $v_0 = n$ vertices. Then
\[
  \prod_{i=0}^{\text{steps}-1} \frac{v_i - 2}{v_i}
  \;=\;
  \prod_{i=0}^{n-3} \frac{n - i - 2}{n - i},
\]
where each factor is taken as a real number, the difference $v_i - 2$ in the left-hand
numerator being the truncated subtraction in $\bbn$ and the right-hand differences being
taken in $\bbr$.
\end{lemma}

\begin{theorem}[Product lower bound for min-cut survival in Karger contraction]
\lean{KargerTrace.survivingTrace_product_survival_lower_bound}
\label{KargerTrace.survivingTrace_product_survival_lower_bound}
\leanok
\uses{IsMulMinCut, KargerTrace.SurvivingTrace, KargerTrace.survivingTrace_factor_product_eq_telescope, KargerTrace.survivingTrace_factor_product_le_survival_product, MulCut.size, Multigraph.edgeCount, Multigraph.vertexCount, survivalProb, telescope_prod}
\difficulty{3}
Let $T$ be an indexed surviving trace of $n-2$ steps over a vertex type $\alpha$, with
states $s_0, s_1, s_2, \ldots$, and suppose $2 \le n$. Write $c$ for the size of the
distinguished cut of the initial state $s_0$, and $m_i$ for the edge count of the
multigraph at state $s_i$. Assume that the distinguished cut of $s_0$ is a minimum cut
of its multigraph, that the multigraph of $s_0$ has exactly $n$ vertices, and that for
every $i < n-2$ the multigraph of $s_i$ has at least two vertices. Then
\[
  \frac{2}{n(n-1)} \;\le\; \prod_{i=0}^{n-3}\left(1 - \frac{c}{m_i}\right),
\]
where each factor $1 - c/m_i$ is the step-survival probability determined by the initial
cut size $c$ and the edge count $m_i$.
\end{theorem}

\begin{definition}[Product of uniform survival probabilities]
\lean{KargerTrace.survivingTraceUniformSurvivalProduct}
\label{KargerTrace.survivingTraceUniformSurvivalProduct}
\leanok
\uses{KargerTrace.SurvivingTrace, KargerTrace.uniformSurvivalProb}
The product over $i < \text{steps}$ of the uniform one-step survival probabilities of the
states of an indexed surviving trace.
\end{definition}

\begin{lemma}[Uniform survival product equals the product of one-step survival probabilities]
\lean{KargerTrace.survivingTrace_uniform_product_eq_survival_product}
\label{KargerTrace.survivingTrace_uniform_product_eq_survival_product}
\leanok
\uses{KargerTrace.SurvivingTrace, KargerTrace.survivingStep_edgeCount_pos, KargerTrace.survivingTraceUniformSurvivalProduct, KargerTrace.uniformSurvivalProb_eq_survivalProb, MulCut.size, Multigraph.edgeCount, survivalProb}
\difficulty{2}
Let $T$ be an indexed surviving trace of length $\text{steps}$, recording at each index
$i$ a state consisting of a multigraph $G_i$ together with a distinguished cut $C_i$,
with consecutive states below $\text{steps}$ related by a surviving step. Then the
product of the uniform one-step survival probabilities of the states of $T$ equals the
product, over all $i < \text{steps}$, of the one-step survival probability $1 -
\dfrac{c_i}{m_i}$, where $c_i$ is the size of the cut $C_i$ (its number of crossing
edges) and $m_i$ is the total edge count of $G_i$; that is,
\[
  \prod_{i < \text{steps}} \frac{\#\{\text{non-crossing edges of state } i\}}{m_i}
  \;=\;
  \prod_{i < \text{steps}} \left(1 - \frac{c_i}{m_i}\right).
\]
\end{lemma}

\begin{lemma}[Uniform survival product via the initial cut size]
\lean{KargerTrace.survivingTrace_uniform_product_eq_initial_cut_survival_product}
\label{KargerTrace.survivingTrace_uniform_product_eq_initial_cut_survival_product}
\leanok
\uses{KargerTrace.SurvivingRun, KargerTrace.SurvivingTrace, KargerTrace.survivingRun_cut_size_eq, KargerTrace.survivingTraceUniformSurvivalProduct, KargerTrace.survivingTrace_prefix_run, KargerTrace.survivingTrace_uniform_product_eq_survival_product, MulCut.size, Multigraph.edgeCount, survivalProb}
\difficulty{3}
Let $T$ be an indexed surviving trace of length $\mathrm{steps}$, with states $s_0, s_1,
\dots$, and write $c$ for the size of the distinguished cut of the initial state $s_0$
and $m_i$ for the edge count of the graph of state $s_i$. Then the product of the
uniform one-step survival probabilities of the states $s_0, \dots, s_{\mathrm{steps}-1}$
equals
\[
  \prod_{i=0}^{\mathrm{steps}-1} \left(1 - \frac{c}{m_i}\right),
\]
the product of the one-step survival probabilities computed with the initial cut size
$c$ at every index.
\end{lemma}

\begin{theorem}[Karger survival lower bound for a minimum cut]
\lean{KargerTrace.survivingTrace_uniform_survival_product_lower_bound}
\label{KargerTrace.survivingTrace_uniform_survival_product_lower_bound}
\leanok
\uses{IsMulMinCut, KargerTrace.SurvivingTrace, KargerTrace.survivingTraceUniformSurvivalProduct, KargerTrace.survivingTrace_product_survival_lower_bound, KargerTrace.survivingTrace_uniform_product_eq_initial_cut_survival_product, Multigraph.vertexCount}
\difficulty{2}
Let $n \ge 2$ and let $T$ be a surviving trace of $n-2$ steps, indexing states $T_0,
T_1, \ldots$ each consisting of a multigraph together with a distinguished cut. Suppose
the distinguished cut of the initial state $T_0$ is a minimum cut of its graph, that
this initial graph has exactly $n$ vertices, and that for every step index $i < n-2$ the
graph of $T_i$ has at least two vertices. Then the product of the uniform one-step
survival probabilities of the states $T_0, \ldots, T_{n-3}$ satisfies \[
\frac{2}{n(n-1)} \le \prod_{i=0}^{n-3} p_i, \] where $p_i$ is the number of edges of
$T_i$ not crossing its distinguished cut divided by the total edge count of $T_i$.
\end{theorem}

\begin{definition}[Product of occurrence-space survival ratios]
\lean{KargerTrace.survivingTraceEdgeChoiceSurvivalProduct}
\label{KargerTrace.survivingTraceEdgeChoiceSurvivalProduct}
\leanok
\uses{KargerTrace.SurvivingTrace, KargerTrace.edgeChoiceSurvivalRatio}
The product over $i < \text{steps}$ of the finite sample-space survival ratios of the
states of an indexed surviving trace.
\end{definition}

\begin{lemma}[Occurrence-space and uniform survival products agree]
\lean{KargerTrace.survivingTrace_edgeChoice_product_eq_uniform_product}
\label{KargerTrace.survivingTrace_edgeChoice_product_eq_uniform_product}
\leanok
\uses{KargerTrace.SurvivingTrace, KargerTrace.edgeChoiceSurvivalRatio_eq_uniformSurvivalProb, KargerTrace.survivingTraceEdgeChoiceSurvivalProduct, KargerTrace.survivingTraceUniformSurvivalProduct}
\difficulty{2}
Let $T$ be an indexed surviving trace of $\text{steps}$ steps, so that $T$ assigns to
each index $i \in \bbn$ a state $T_i$ and, for every $i < \text{steps}$, the passage
from $T_i$ to $T_{i+1}$ is a surviving step. Then
\[
\prod_{i \,<\, \text{steps}} r(T_i) \;=\; \prod_{i \,<\, \text{steps}} p(T_i),
\]
where $r(T_i)$ denotes the occurrence-space survival ratio of $T_i$ (the number of
surviving edge occurrences divided by the total number of edge occurrences) and $p(T_i)$
denotes its uniform one-step survival probability (the number of non-crossing edges
divided by the total edge count).
\end{lemma}

\begin{theorem}[Karger survival product lower bound]
\lean{KargerTrace.survivingTrace_edgeChoice_survival_product_lower_bound}
\label{KargerTrace.survivingTrace_edgeChoice_survival_product_lower_bound}
\leanok
\uses{IsMulMinCut, KargerTrace.SurvivingTrace, KargerTrace.survivingTraceEdgeChoiceSurvivalProduct, KargerTrace.survivingTrace_edgeChoice_product_eq_uniform_product, KargerTrace.survivingTrace_uniform_survival_product_lower_bound, Multigraph.vertexCount}
\difficulty{1}
Fix an integer $n \ge 2$ and let $T$ be an indexed surviving trace of length $n-2$: a
sequence of states $T_0, T_1, \dots$ in which each $T_{i+1}$ (for $i < n-2$) arises from
$T_i$ by contracting an edge of its multigraph that does not cross the distinguished
cut. Suppose the distinguished cut of the initial state $T_0$ is a minimum cut of the
initial multigraph, that this initial multigraph has exactly $n$ vertices, and that the
multigraph of $T_i$ has at least $2$ vertices for every $i < n-2$. Then the product over
$i < n-2$ of the one-step survival ratios $\mathrm{surv}(T_i)$ — each the fraction of
edge occurrences of $T_i$ whose underlying edge does not cross the distinguished cut —
satisfies \[ \frac{2}{n(n-1)} \le \prod_{i=0}^{n-3} \mathrm{surv}(T_i). \]
\end{theorem}
